\scanpage{317}\mspage{41}
[Nous l'appliquerons dans les cas suivants
\[
\mathcal F=\mathcal F_{W^{\deleted{\unsure{k|*}}}},\mathrm{Fib}f_{W^{\deleted{\unsure{k|*}}}},\mathrm{Perf}_W,
\deleted{\unsure{\mathrm{Perf}_{W_0}}}\mathcal F_{W_0^{\deleted{\unsure{k|*}}}},
\mathrm{Fib}f_{W_0^{\unsure{k|*}}},\mathrm{Perf}_{W_0}.
\tag*{(4.4.1)}
\]
Je signale d'ailleurs que l'on verra plus bas \inserted{$\underbrace{\text{(cor. }\unsure{4.4.10}\text{)}}$} que
\[
 A(\mathcal F_W)\Longleftrightarrow A(\mathrm{Fib}f_W),\quad
 A(\mathcal F_{W_0^{\unsure{k|*}}})\Longleftrightarrow A(\mathrm{Fib}f_{W_0^{\unsure{k|\alpha|*}}}),
 \tag*{(4.4.3)}
\]
et de même pour les variantes $B,C$, de
sorte que pour l'application que nous voulons
étudier, il n'y a pas lieu de distinguer les
$\mathcal F_W$ et $\mathrm{Fib}f_W$, les $\mathcal F_{W_0}$ et les $\mathrm{Fib}f_{W_0}$.
D'ailleurs, il suffira de regarder les \deleted{\unsure{catég.}}
ensembles $\mathcal F_{W_0},\mathrm{Fib}f_{W_0},\mathrm{Perf}_{W_0}$, puisque \unsure{le|les}
cas \deleted{des} \deleted{\unsure{catég.}} \unsure{relatif|relatifs} à $W$ en sont un cas
particulier (en prenant $W_0=W$).]
% slanted
\margin{Cela fera donc a priori $3\times6=18$ conditions différentes, mais à cause de (4.4.3) il n'y en a (à \unsure{équiv.|isom.|ignor.} près) que $4\times3=12$, et en fait deux parmi ces quatre triplets de conditions se réduisent à une seule (pour $\mathcal F=\mathcal F_W$ et $\mathcal F=\mathcal F_{W_0}$), de sorte qu'il ne reste au total que \uline{huit} conditions différentes, que nous \unsure{résumons|désignons} dans le diagramme récapitulatif (\unsure{4.5.1|4.5.1.1}) de la page \unsure{48|40}. Il y en a \deleted{\unsure{en}} seulement \uline{quatre} $A_{W_0},B_{W_0},C_{W_0},D_{W_0}$ \unsure{comme|à donner comme} fonctions de $W_0$, les quatre autres $A,B,C,D$ s'obtiennent en \unsure{en} spécialisant au cas où $W_0=W$.}

\deleted{Cela résultera évidemment du}

\deleted{\uline{Lemme} 4.4.1 Soit $f:X\text{---}Y$ un $S$-morphisme dans $\mathrm{Cat}$, avec $p:X\text{---}S$ et $q:Y\longrightarrow S$ des $W$-fibrations. Si $f\in W$, donc $f\in W^u$. Cela résulte en effet formellement du fait que $\mathrm{Cat}$ a assez de $W$-fibrations}
% slanted
\margin{Vaut dans une catégorie de modèles quelconque.}

\uline{Lemme} 4.4.1. Supposons que $\mathcal F$ satisfasse
la condition suivante : toute flèche $f$ de
$\mathrm{Cat}$ se factorise en $f=pi$, avec $i\in W$, $f\in\mathcal F$.
Alors la condition $A(\mathcal F)$ sur $\mathcal F$ implique
que dans $\mathrm{Cat}$ il y a ``assez de $W$-fibrations'',
et par suite (par un argument formel bien
% archive sheet number: 316
